% Part: computability
% Chapter: computability-theory
% Section: coding-computations

\documentclass[../../../include/open-logic-section]{subfiles}

\begin{document}

\olfileid{cmp}{thy}{cod}
\olsection{编码计算}

在每一种计算模型中，都可以进行以下操作：

\begin{enumerate}
\item 以系统化的方式描述可计算函数的\emph{定义}。例如，Turing机规格、递归定义或编程语言中的程序，都可以看作提供了这些定义。
\item 描述由某个定义给出的函数在给定输入上的计算过程的完整记录。例如，一次Turing机计算可以通过每一步计算的格局（机器状态、纸带内容）序列来描述。
\item 检验一个假定的计算记录，是否确实是具有给定定义的可计算函数在给定输入上的计算记录（函数在该输入上有定义，即计算停机）。
\item 从对计算过程完整记录的这类描述中，提取函数在给定输入上的值。例如，一次停机的Turing机计算中，最后一步的纸带内容即为该值。
\end{enumerate}

利用编码，可以为每个可计算函数的描述分配一个数值\emph{指标}，使得我们能够以可计算的方式从该指标恢复出相应的指令。类似地，一次计算过程的完整记录也能用单个数来编码。由此得到的算术关系“$s$~编码了指标为~$e$ 的函数在输入~$x$ 上的计算记录”以及函数“代码为~$s$ 的计算序列的输出”都是可计算的；事实上，它们是原始递归的。

这一基本事实威力巨大，它使得我们能够独立于所选择的计算模型，证明一系列关于可计算性的深刻而重要的结果。

\end{document}